% Figure: righting-arm curve GZ(theta) for a floating body. Rises from the
% origin with initial slope GM (radians), peaks, then falls to zero at the angle
% of vanishing stability. The initial tangent line of slope GM is also drawn.
% No em-dashes.
\begin{figure}[htbp]
\centering
\begin{tikzpicture}[scale=1.0]
\begin{axis}[
  width=11cm, height=5.6cm,
  xlabel={heel angle $\theta$ (degrees)},
  ylabel={righting arm $\overline{GZ}$ (m)},
  xmin=0, xmax=80, ymin=0, ymax=0.85,
  xtick={0,15,30,45,60,75},
  ytick={0,0.25,0.5,0.75},
  tick label style={font=\scriptsize},
  label style={font=\scriptsize},
  legend pos=north east,
  legend style={font=\scriptsize, draw=enggray},
  axis lines=left,
  samples=120, domain=0:72,
]
% GZ curve: a smooth hump that returns to zero near 72 deg.
% Use GZ = A*sin(theta)*(1 - theta/theta_v) shape, clamped non-negative by domain.
\addplot[engblue, very thick] ({x}, {0.95*sin(x)*(1 - x/72)});
\addlegendentry{$\overline{GZ}(\theta)$}
% initial tangent of slope GM per radian: GZ ~ GM*theta. GM chosen 0.55 m.
% slope in (m per degree) = GM * pi/180 = 0.55*0.01745 = 0.0096
\addplot[engamber, thick, dashed, domain=0:57] ({x}, {0.0096*x});
\addlegendentry{initial slope $=\overline{GM}$}
% mark the peak (max righting arm) and the angle of vanishing stability
\addplot[only marks, engred, mark=*, mark size=1.6pt] coordinates {(72,0)};
\node[engred, font=\scriptsize, anchor=south west] at (axis cs:60,0.04) {angle of vanishing stability};
\end{axis}
\end{tikzpicture}
\caption[The righting-arm curve]{The righting-arm curve
$\overline{GZ}(\theta)$ for a floating body. Near upright the curve rises
with initial slope equal to the metacentric height $\overline{GM}$ (the
amber tangent), which is why $\overline{GM}$ governs small-angle stability.
At larger heel the geometry of the immersed wedge changes, the curve peaks
(the maximum righting arm sets the largest steady heeling moment the body
can resist), and it then falls to zero at the angle of vanishing stability
(red), beyond which the couple becomes capsizing. The area under the curve
up to a given angle is the dynamic stability, the energy the body can
absorb from a gust or wave before capsize.}
\label{fig:vol03ch10:righting-arm}
\end{figure}
